\begin{abstract}
This project predicts the daily bike rental demand for the 20 most used Capital Bikeshare start stations. 
We frame this as a supervised regression problem with \code{total\_rentals} as the target variable. 
The main focus is not just on the final model score, but on showing a clear and reproducible process 
from raw data to the final evaluation. We combine raw trip, station, and weather data into daily 
records for each station. To prepare the data, we remove variables that could cause data leakage, 
drop redundant weather features using correlation analysis, apply one-hot encoding for station IDs, 
and add past rental data (lag features). All models are evaluated using a chronological 
train-validation-test split to respect the timeline of the data. Our best model is Gradient Boosting 
with lag features, which achieves a test RMSE of 22.78, a MAE of 16.92, and an $R^2$ of 0.663. 
The clearest result of our project is that time factors matter the most: adding lag features 
improved the test RMSE for every model we evaluated.
\end{abstract}